\documentclass[11pt,a4paper]{article}
\usepackage[utf8]{inputenc}
\usepackage[T1]{fontenc}
\usepackage[brazil]{babel}
\usepackage[margin=2.3cm]{geometry}
\usepackage{booktabs}
\usepackage{amssymb}
\usepackage{array}
\usepackage{enumitem}
\usepackage{xcolor}
\usepackage{titlesec}
\usepackage{parskip}
\usepackage[hidelinks]{hyperref}
\usepackage{fancyhdr}

\definecolor{puc}{RGB}{0,60,113}
\titleformat{\section}{\large\bfseries\color{puc}}{\thesection}{0.6em}{}
\titleformat{\subsection}{\normalsize\bfseries}{\thesubsection}{0.6em}{}
\setlist{nosep,leftmargin=1.4em}

\pagestyle{fancy}\fancyhf{}
\lhead{\small\color{puc}Avanços — Caso de Replicação Heine (Xande)}
\rhead{\small\thepage}
\renewcommand{\headrulewidth}{0.4pt}

\begin{document}

\begin{center}
{\LARGE\bfseries\color{puc} Avanços do Caso de Replicação: Heine (Xande)}\\[2pt]
{\large Amostragem em Grandes Volumes de Dados em Redes Sociais Digitais}\\[6pt]
{\normalsize Dissertação de Mestrado — Mariana Barreto \textbullet\ PUC-Rio, Informática}\\
{\small Orientador: Prof.\ Sérgio Lifschitz (grupo eTC/TriTech) \textbullet\ Relatório de progresso — 17/07/2026}
\end{center}

\vspace{2pt}
\noindent\fcolorbox{puc}{puc!5}{\parbox{\dimexpr\textwidth-2\fboxsep-2\fboxrule}{%
\small\textbf{Em uma frase.} Este caso replica, por \emph{expansão}, a dissertação de
mestrado de Alexandre Heine (2025) — o precursor metodológico direto da linhagem do
grupo — aplicando a ela o mesmo fluxo já usado no caso Twitter/Ituassu. O trabalho até
aqui deixou \textbf{a documentação e o \emph{pipeline} de análise prontos} e concluiu a
\textbf{Fase~0 (descoberta dos dados)}, que revelou onde a coleta original está — e onde
não está. A execução das análises aguarda o \emph{download} dos dados (SharePoint do BioBD).}}

\section{O caso e por que ele importa}

A dissertação de Heine \& Lifschitz (2025), \emph{Um Estudo sobre Amostragem em Grandes
Volumes de Dados em RSD}, comparou \textbf{amostra estatística $\times$ conjunto completo}
em $\sim$57,9 milhões de \emph{tweets} das eleições brasileiras de 2022 (Lula $\times$
Bolsonaro), em duas análises: \textbf{modelagem de tópicos} (BERTopic) e \textbf{estatística
quantitativa} (engajamento e distribuição temporal, via SGBD).

O encaixe na dissertação da Mariana é privilegiado: enquanto os demais casos do
\emph{lineup} tratam autores que \emph{sub-coletaram sem perceber}, \textbf{o próprio Heine
já mediu os dois lados} — e concluiu que \emph{a amostra falha} em reproduzir os tópicos do
todo, \emph{mesmo usando o tamanho de amostra calculado pela fórmula de proporção}
(99\,\% de confiança, $\pm$1\,\%). Mais: nos Trabalhos Futuros ele \emph{deixa em aberto} a
suspeita de que a fórmula de proporção seria inadequada para uma ``medida composta'' como o
sentido de um texto. \textbf{É essa fresta que a replicação ocupa.}

\section{O fluxo seguido (espelhando o caso Twitter/Ituassu)}

Reproduzimos, para o caso Heine, a mesma estrutura do caso já avançado do grupo:

\begin{itemize}
\item \textbf{Documentação da replicação:} plano de fases, ficha do alvo e um documento por
eixo de análise (tópicos e quantitativo), com \emph{predições pré-registradas}.
\item \textbf{Infraestrutura de conexão:} acesso ao banco da PUC via túnel SSH sobre a VPN
Cloud-DI, reaproveitando o módulo do caso Twitter.
\item \textbf{\emph{Pipeline} reprodutível:} diagnóstico do banco (Fase~0) $\rightarrow$
\emph{snapshot} congelado local (Fase~1) $\rightarrow$ curva de convergência e análises
(Fases~2--3) $\rightarrow$ tabela mestre ``mudou? / convergiu?'' (Fase~4).
\end{itemize}

Os hiperparâmetros do BERTopic foram \textbf{extraídos do repositório do próprio Heine}
(\texttt{Maxteralex/representative-dataset-problem}) para garantir fidelidade:
\emph{embeddings} \texttt{paraphrase-multilingual-MiniLM-L12-v2}; UMAP
($n\_neighbors{=}15$, $n\_components{=}5$, $min\_dist{=}0$, cosseno, \texttt{random\_state=42});
HDBSCAN ($min\_cluster\_size{=}15$); e vetorizador com \emph{stopwords} em português.

\section{Descoberta dos dados (Fase 0) — concluída}

A Fase~0 foi uma investigação sistemática para localizar a coleta de 2022. Foram varridos
\textbf{3 servidores (vm031, vm037, vm067), 29 bancos PostgreSQL, 3 MongoDBs}, o disco
local e o repositório do Heine no GitHub. Resultado:

\begin{center}\small
\begin{tabular}{@{}p{5.2cm}p{9.3cm}@{}}
\toprule
\textbf{Onde} & \textbf{O que tem} \\
\midrule
Postgres \texttt{eTC\_Producao} (vm067) & Coleta 2014 (\texttt{ituassu\_2014}, caso Twitter) e 2018; 2022 apenas testes. \\
Postgres \texttt{eTC\_Producao} (vm031) & Coleta contínua; a busca eleitoral 2022 tem só $\sim$25 mil \emph{tweets} (o dia 2/out tem $\sim$8 mil — vs.\ 1,29 milhão do Heine). \\
vm037 & Servidor das \emph{disciplinas} de Banco de Dados (bancos de alunos). \\
MongoDB (vm031) & \textbf{Coleta \emph{Instagram} do Heine} (\texttt{xande\_search}, $\sim$628 mil posts) \textbf{+ Reddit}. \\
\bottomrule
\end{tabular}
\end{center}

\textbf{Achado central:} o \emph{Twitter} de 57,9\,M / 1,29\,M-do-dia-2 que a dissertação
usa \textbf{não está nos bancos acessíveis} — provavelmente foi arquivado após a defesa
(abr/2025). O que está acessível é a coleta \emph{Instagram} do Heine (628\,k) e uma coleta
\emph{Reddit} (que servem, de bônus, a outros casos do \emph{lineup}). O dado de Twitter
está no \textbf{SharePoint do BioBD PUC-Rio}, acessível pela conta institucional da Mariana.
O repositório do Heine mostrou ainda que as análises dele leem de \textbf{CSVs}
(\texttt{1st\_turn\_day.csv}), coluna \texttt{content\_text}.

\section{O que já roda quando o dado chegar}

O \emph{pipeline} está pronto e aceita tanto \textbf{CSV} (coluna \texttt{content\_text},
como o Heine usou) quanto \emph{snapshot} local:

\begin{itemize}
\item \textbf{Eixo tópicos:} treina o BERTopic no dia completo e mede, para \emph{frações
crescentes} do volume (com réplicas \emph{bootstrap}), a distância entre a distribuição de
tópicos da fração e a do todo (divergência de Jensen--Shannon), sobrepondo o $n$ da fórmula
de proporção do Heine para mostrar \emph{onde a estrutura de tópicos realmente estabiliza}.
\item \textbf{Eixo quantitativo:} recalcula engajamento por faixa e a série temporal
diretamente no servidor e mede em que fração de dados essas conclusões já convergem.
\end{itemize}

\section{Perguntas e predições pré-registradas (o valor científico)}

\begin{itemize}
\item \textbf{Contraste-chave:} a hipótese é que o \emph{quantitativo converge cedo}
(coletar mais é barato \emph{e} dispensável) enquanto os \emph{tópicos convergem tarde ou
nunca} dentro do orçamento da fórmula. O mesmo objeto (``coletar mais'') com respostas
opostas \emph{por tipo de análise} — exatamente a tese da dissertação.
\item \textbf{Correção do diagnóstico do Heine:} a fórmula $n_0 = Z^2pq/e^2$ governa a
estabilidade de \emph{uma proporção escalar}, não de um objeto de alta dimensão (a
distribuição de tópicos). A contribuição é \emph{medir a curva real} e mostrar
empiricamente que a estabilização exige muito mais do que o $n$ da fórmula — respondendo à
pergunta que o próprio Heine deixou em aberto.
\end{itemize}

\section{Estado atual e próximos passos}

\begin{center}\small
\begin{tabular}{@{}p{10.8cm}c@{}}
\toprule
\textbf{Item} & \textbf{Estado} \\
\midrule
Documentação da replicação (plano, alvo, eixos) & \checkmark\ pronto \\
\emph{Pipeline} (diagnóstico, \emph{snapshot}, convergência, quantitativo) & \checkmark\ pronto \\
Fase 0 — descoberta e caracterização dos dados & \checkmark\ concluída \\
Organização em \texttt{Documents/dissertacao/} (escrita + casos) & \checkmark\ feita \\
Fases 1--4 — execução das análises & aguardando dados \\
\bottomrule
\end{tabular}
\end{center}

\textbf{Falta apenas} baixar do SharePoint do BioBD os CSVs (\texttt{1st\_turn\_day.csv} e
\texttt{1st\_turn\_day\_random\_sample.csv}) ou o \emph{dump} do MongoDB de Twitter. Com o
arquivo em mãos, as Fases~1--4 rodam de imediato e produzem a curva de convergência, as
tabelas comparativas e a seção correspondente do texto da dissertação — sempre com números
rastreáveis a \emph{scripts} e consultas.

\vspace{4pt}
\noindent\rule{\textwidth}{0.4pt}\\[-2pt]
{\footnotesize Documentos de apoio (na pasta do caso): \texttt{REPLICACAO\_CASO\_HEINE.md},
\texttt{DISSERTACAO\_HEINE\_alvo.md}, \texttt{EIXO\_TOPICOS.md},
\texttt{EIXO\_QUANTITATIVO.md} e o \emph{log} de descoberta
\texttt{data/repl/heine2022/FASE0\_descoberta.md}.}

\end{document}
